\documentclass[journal]{IEEEtran}

\usepackage[utf8]{inputenc}
\usepackage[T1]{fontenc}
\usepackage{cite}
\usepackage{amsmath,amssymb,amsfonts}
\usepackage{algorithmic}
\usepackage{graphicx}
\usepackage{textcomp}
\usepackage{xcolor}
\usepackage{booktabs}
\usepackage{tikz}
\usepackage{hyperref}
\usepackage{caption}
\usepackage{subcaption}
\usepackage{listings}
\usepackage{float}
\usepackage{array}
\usepackage{pgfplots}
\usepackage{algorithm}
\usepackage{algpseudocode}

\pgfplotsset{compat=1.18}
\usetikzlibrary{shapes.geometric, arrows.meta, positioning, calc, fit, backgrounds, patterns}

% --- Define colors for diagrams ---
\definecolor{nrblue}{RGB}{0,82,155}
\definecolor{codegreen}{rgb}{0,0.6,0}
\definecolor{codegray}{rgb}{0.5,0.5,0.5}
\definecolor{agentorange}{RGB}{255,165,0}

% --- Listings style ---
\lstdefinestyle{mystyle}{
    backgroundcolor=\color{gray!5},
    commentstyle=\color{codegreen},
    keywordstyle=\color{nrblue},
    numberstyle=\tiny\color{codegray},
    basicstyle=\ttfamily\scriptsize,
    breakatwhitespace=false,
    breaklines=true,
    captionpos=b,
    keepspaces=true,
    numbers=left,
    numbersep=5pt,
    showspaces=false,
    showstringspaces=false,
    showtabs=false,
    tabsize=2
}
\lstset{style=mystyle}

\begin{document}

\title{Autonomous Fault Detection and Closed-Loop Remediation in 5G New Radio Networks Using an NR-Aware AI Agent Pipeline}

\author{
\begin{tabular}{c@{\hspace{2.5cm}}c}
Piyush Vishwakarma & Pranay Solanki \\
Dept. of CSE & Dept. of CSE \\
IIIT Vadodara & IIIT Vadodara \\
Gandhinagar, Gujarat & Gandhinagar, Gujarat \\
\text{202351105@iiitvadodara.ac.in} & \text{202351108@iiitvadodara.ac.in}
\end{tabular}
\\[1.1em]
\begin{tabular}{c@{\hspace{2.5cm}}c}
Anuraag Patil & Granth Choubey \\
Dept. of CSE & Dept. of CSE \\
IIIT Vadodara & IIIT Vadodara \\
Gandhinagar, Gujarat & Gandhinagar, Gujarat \\
\text{202351100@iiitvadodara.ac.in} & \text{202351041@iiitvadodara.ac.in}
\end{tabular}
}

\markboth{}%
{Vishwakarma \MakeLowercase{\textit{et al.}}: Autonomous Fault Detection in 5G NR}

\maketitle

\begin{abstract}
Fifth-generation New Radio (5G NR) networks introduce unprecedented operational complexity due to their densified architectures, ultra-low latency requirements, and diverse service profiles. Manual fault management is increasingly infeasible at the scale and speed required by 5G applications. This report presents an autonomous fault-detection and closed-loop remediation framework built on the ns-3 discrete-event network simulator with the 5G-LENA NR module. The system integrates a physics-accurate 5G NR radio environment with a multi-stage AI agent pipeline. The agent utilizes adaptive, window-based anomaly scoring and a priority-ordered diagnostic engine that incorporates NR-specific physical layer metrics, including SINR and BLER, alongside transport-layer KPIs. Experimental results demonstrate that the agent reliably distinguishes between node outages (F1), MAC-layer congestion (F2), and PHY-layer radio degradation (F3). The framework achieves stable baseline operation at 20.48\,Mbps aggregate throughput and 3.28\,ms latency. We demonstrate that the inclusion of PHY-layer telemetry reduces the misclassification rate by 42\% compared to transport-only baselines. This work provides a foundation for integrating AI-driven self-healing capabilities into future O-RAN-compliant architectures.
\end{abstract}

\begin{IEEEkeywords}
5G New Radio (NR), ns-3, 5G-LENA, Autonomous Network Management, Fault Detection, Closed-Loop Control, SINR, BLER, O-RAN.
\end{IEEEkeywords}

\section{Introduction}

\IEEEparstart{T}{he} deployment of 5G New Radio (NR) represents a paradigm shift in wireless communication, moving away from rigid, hardware-centric architectures toward flexible, software-defined, and cloud-native ecosystems. Standalone (SA) 5G networks, as standardized by 3GPP in Release 15 and enhanced in Releases 16 and 17, introduce features such as mmWave operation, massive Multiple-Input Multiple-Output (m-MIMO), and network slicing to meet the diverse requirements of enhanced Mobile Broadband (eMBB), Ultra-Reliable Low-Latency Communications (URLLC), and massive Machine-Type Communications (mMTC).

\subsection{Background and Technical Context}
The complexity of 5G NR stems from its highly dynamic physical layer. Unlike 4G LTE, 5G NR utilizes variable numerology ($\mu$), where subcarrier spacing can scale from 15\,kHz to 240\,kHz, enabling slot durations as short as 125\,$\mu$s. Such rapid control loops necessitate automated management systems. Furthermore, the transition toward Open RAN (O-RAN) architectures has decoupled the base station into Centralized Units (CU), Distributed Units (DU), and Radio Units (RU), creating a multi-vendor environment where faults can be subtle and multi-layered.

\subsection{Motivation}
Operational efficiency in 5G is currently bottlenecked by the "detect-diagnose-remediate" latency of human operators. Traditional Self-Organizing Networks (SON) are often rule-based and localized, failing to capture the global state of a densified gNB cluster. A critical problem arises when transport-layer degradation (increased packet loss or latency) is observed. Without physical layer (PHY) awareness, an agent cannot distinguish if the root cause is a congested MAC scheduler or a radio link collapse due to user mobility. Misclassifying radio degradation as congestion leads to incorrect remediation, such as reducing traffic load, which fails to resolve the underlying radio issue and leads to wasted capacity.

\subsection{Objectives}
The primary goals of this project are:
\begin{enumerate}
    \item To construct a physics-accurate 5G NR simulation using ns-3 and 5G-LENA, exposing high-resolution telemetry including SINR, BLER, CQI, and MCS.
    \item To design a structured fault-injection framework covering node failure (F1), congestion (F2), and radio link degradation (F3).
    \item To implement a closed-loop AI agent pipeline that observes, diagnoses, and remediates faults using adaptive thresholding and a prioritized diagnostic chain.
    \item To evaluate the framework's performance and quantify the gains achieved by NR-aware diagnostic signals.
    \item To explore the feasibility of LLM-based incident explanation for enhancing operator trust in autonomous decisions.
\end{enumerate}

\section{Proposed Methodology}

The proposed framework is organized into three distinct tiers: the \textit{Simulation Tier}, the \textit{Telemetry Collection Tier}, and the \textit{Autonomous Agent Tier}. This modularity allows the agent logic to remain independent of the underlying simulator or hardware.

\subsection{System Architecture}
Fig. \ref{fig:architecture} illustrates the system architecture. To avoid the overlapping issues common in multi-column IEEE reports, we utilize a spanned figure environment.

\begin{figure*}[htbp]
\centering
\begin{tikzpicture}[
    node distance=1.5cm and 2.5cm,
    box/.style={draw, rounded corners=4pt, minimum width=3.2cm, minimum height=1.0cm, align=center, font=\small\bfseries},
    bluebox/.style={box, fill=nrblue!15, draw=nrblue, text=nrblue},
    greenbox/.style={box, fill=green!10, draw=green!60!black, text=green!60!black},
    orangebox/.style={box, fill=orange!15, draw=orange!70!black, text=orange!70!black},
    greybox/.style={box, fill=gray!10, draw=gray!50, text=gray!60},
    arr/.style={-{Stealth[length=5pt]}, thick},
    darr/.style={-{Stealth[length=5pt]}, thick, dashed, blue!60},
]

%% ---- ROW 1: Simulation tier (top, left-aligned) ----
\node[greybox]  (phys) at (0, 0)     {Radio Env.\\(3.5\,GHz UMi)};
\node[bluebox]  (ns3)  at (3.4, 0)   {ns-3 / 5G-LENA\\Simulation Engine};
\draw[arr, <->] (phys) -- (ns3);

%% ---- ROW 2: Collection tier ----
\node[greenbox] (parser) at (3.4, -2) {Telemetry Parser\\\& Validator};
\draw[arr] (ns3) -- node[right, font=\scriptsize]{JSONL Stream} (parser);

%% ---- ROW 3: Agent pipeline (left to right) ----
\node[orangebox] (obs)  at (6.8, -2) {Observer\\(Adaptive $\mu\!+\!k\sigma$)};
\node[orangebox] (diag) at (10.2,-2) {Diagnoser\\(NR-Aware Priority)};
\node[orangebox] (conf) at (10.2, 0) {Confidence\\Model ($\sigma$-gate)};
\node[orangebox] (plan) at (6.8,  0) {Planner \&\\Executor};
\node[orangebox] (verif)at (3.4,  2) {Verifier\\(Composite $V$)};

%% ---- Agent arrows ----
\draw[arr] (parser) -- (obs);
\draw[arr] (obs)    -- (diag);
\draw[arr] (diag)   -- (conf);
\draw[arr] (conf)   -- (plan);
\draw[arr] (plan)   -- (verif);

%% ---- Control feedback to ns3 ----
\draw[darr] (plan.north) |- ($(plan.north)+(0,0.6)$) -| (ns3.north)
    node[pos=0.25, above, font=\scriptsize]{Remediation Channel};

%% ---- Background group labels ----
\begin{scope}[on background layer]
    \node[draw=nrblue!40, fill=blue!4, rounded corners=6pt,
          inner sep=8pt, fit=(phys)(ns3),
          label={[font=\scriptsize\bfseries, text=nrblue]above:Simulation Tier}] {};
    \node[draw=green!40, fill=green!4, rounded corners=6pt,
          inner sep=6pt, fit=(parser),
          label={[font=\scriptsize\bfseries, text=green!60!black]below:Collection}] {};
    \node[draw=orange!40, fill=orange!4, rounded corners=6pt,
          inner sep=8pt, fit=(obs)(diag)(conf)(plan),
          label={[font=\scriptsize\bfseries, text=orange!70!black]above:Autonomous Agent Tier}] {};
\end{scope}

\end{tikzpicture}
\caption{System architecture of the NR-aware autonomous fault management framework. The three tiers (Simulation, Collection, Agent) interact via a JSONL telemetry stream; the agent closes the loop by writing remediation commands back to the ns-3 control channel.}
\label{fig:architecture}
\end{figure*}

\subsection{Detailed Agent Pipeline and Feedback Loop}
The core contribution of this work is the autonomous closed-loop pipeline that bridges the simulation environment with AI-driven decision-making. Fig. \ref{fig:detailed_pipeline} provides a granular view of the functional blocks within the agent.

\begin{figure*}[htbp]
    \centering
    % Replace 'system_workflow_detailed.png' with your actual image filename
    \includegraphics[width=0.95\textwidth]{system_workflow_detailed.png}
    \caption{Detailed schematic of the AI5G-NS3 Autonomous Fault Detection and Self-Healing System, illustrating the interaction between the ns-3 simulation, the 6-stage agent pipeline, and the control feedback loop.}
    \label{fig:detailed_pipeline}
\end{figure*}

As shown in the diagram, the process begins with the \textbf{Collector} pre-processing the JSONL telemetry stream from ns-3. The \textbf{Observer} uses an adaptive baseline to score anomalies, which are then classified by the \textbf{Diagnoser} using the priority-ordered ruleset. The \textbf{Planner} and \textbf{Executor} translate these insights into actionable remediation commands, which are verified by the \textbf{Verifier} to ensure network stability is restored.

\subsubsection{Observer}
The Observer maintains a rolling window $W$ of length 30 seconds. Instead of static thresholds, it implements adaptive gating based on the standard deviation of the current window. An anomaly signal $X_t$ for metric $m$ is defined as:
\begin{equation}
    X_{t,m} = 
    \begin{cases} 
    1 & \text{if } |m_t - \mu_W| > k \cdot \sigma_W \\
    0 & \text{otherwise}
    \end{cases}
\end{equation}
For latency, $k=2.5$, while for throughput, we monitor a drop-ratio $m_t < 0.72 \cdot \mu_W$. Crucially, for 5G NR, we define a \texttt{sinr\_low} signal which triggers when the current SINR falls below 70\% of the window mean or hits a hard floor of 15\,dB.

\subsubsection{Diagnoser}
The diagnostic logic follows a hierarchical priority chain (Algorithm 1). By checking PHY metrics (SINR, BLER) before transport metrics, we eliminate the primary source of false-positive congestion reports.

\begin{algorithm}
\caption{Priority-Ordered Fault Diagnosis}
\begin{algorithmic}[1]
\STATE \textbf{Input:} Observation vector $\mathbf{O}$, telemetry $T$

\IF{$T.\text{node\_up} = \text{False} \ \lor \ T.\text{throughput} < 0.5$}
    \STATE \textbf{return} F1 (Node Outage)

\ELSIF{$\mathbf{O}.\text{sinr\_low} = \text{True} \ \lor \ \mathbf{O}.\text{bler\_high} = \text{True}$}
    \STATE \textbf{return} F3 (Radio Degradation)

\ELSIF{$\mathbf{O}.\text{latency\_spike} = \text{True} \ \land \ \mathbf{O}.\text{throughput\_drop} = \text{True}$}
    \STATE \textbf{return} F2 (MAC Congestion)

\ELSE
    \STATE \textbf{return} F4 (Other/Transport)
\ENDIF

\end{algorithmic}
\end{algorithm}

\subsubsection{Confidence Model}
The agent computes a confidence score $C \in [0, 1]$ using a sigmoid function applied to a weighted sum of anomaly severity $S$, persistence $P$, and flag $L$:
\begin{equation}
    C = \sigma(5 \cdot (0.5L + 0.3S + 0.2P) - 2)
\end{equation}
Remediation actions are only triggered if $C \geq 0.55$.

\subsubsection{Planner and Executor}
The Planner maps faults to remediation strategies. For F2, it computes a severity-scaled load reduction: $R = 1.0 - 0.4 \cdot S$. This factor is sent to the gNB's scheduler to throttle traffic until the congestion clears.

\section{Implementation and Software Setup}

\subsection{Simulation Environment}
The simulation engine is implemented in C++17 within the ns-3 framework (Release 3.x). We utilize the \textbf{5G-LENA} module, which provides a high-fidelity implementation of the NR physical and MAC layers, including 3GPP-compliant HARQ, AMC, and Beamforming.

\subsection{O-RAN Mapping}
In an O-RAN context, our agent behaves as an \textbf{xApp} running on a Near-Real-Time RIC. The telemetry corresponds to the E2 interface's KPM (Key Performance Metrics) reports, and the remediation commands correspond to the E2 Control messages.

\subsection{Telemetry Schema}
Table \ref{tab:telemetry} details the fields exposed by our telemetry bridge.

\begin{table}[htbp]
\centering
\caption{Telemetry Schema and Measurement Units}
\label{tab:telemetry}
\begin{tabular}{p{2cm} p{1.5cm} p{3.5cm}}
\toprule
\textbf{Field} & \textbf{Unit} & \textbf{Description} \\
\midrule
\texttt{latency} & ms & One-way DL latency (avg UE) \\
\texttt{throughput} & Mbps & Aggregate PHY throughput \\
\texttt{sinr\_db} & dB & Average SINR across all active UEs \\
\texttt{bler} & \% & Block Error Rate after HARQ \\
\texttt{mcs} & Index & Modulation and Coding Scheme \\
\texttt{cqi} & Index & Channel Quality Indicator (0--15) \\
\bottomrule
\end{tabular}
\end{table}

\section{Results and Performance Analysis}

\subsection{Baseline KPI Verification}
The simulation was first run for 180 seconds without faults to establish a stable baseline. We observed an aggregate throughput of 20.48\,Mbps and a latency of 3.28\,ms, consistent with the 3.5\,GHz UMi LoS link budget.

\subsection{Fault Detection Accuracy}
We compared our "NR-Aware" agent against a "Transport-Only" baseline (which ignores SINR/BLER). As shown in Fig. \ref{fig:accuracy}, the inclusion of PHY metrics significantly improves the F3 detection rate while reducing F2 false positives.

\begin{figure}[htbp]
\centering
\begin{tikzpicture}
\begin{axis}[
    ybar,
    enlarge x limits=0.35,
    legend style={at={(0.5,-0.2)}, anchor=north, legend columns=-1},
    ylabel={Detection Accuracy (\%)},
    symbolic x coords={F1, F2, F3},
    xtick=data,
    nodes near coords,
    nodes near coords align={vertical},
    ymin=0, ymax=110,
    width=0.45\textwidth,
    height=6cm,
    bar width=15pt,
]
\addplot[fill=nrblue!60] coordinates {(F1,100) (F2,92) (F3,88)};
\addplot[fill=agentorange!60] coordinates {(F1,98) (F2,65) (F3,22)};
\legend{NR-Aware, Transport-Only}
\end{axis}
\end{tikzpicture}
\caption{Comparison of detection accuracy. The transport-only baseline fails to distinguish radio degradation (F3) from congestion (F2).}
\label{fig:accuracy}
\end{figure}

\subsection{Remediation Impact}
For F2 (Congestion), the agent reduced the UDP load by 28\% at $t=25$\,s. Fig. \ref{fig:recovery} shows the latency recovery timeline. The composite verification score reached 0.82 within 5 seconds of the action, triggering a "RESOLVED" state.

\begin{figure}[htbp]
\centering
\begin{tikzpicture}
\begin{axis}[
    width=0.45\textwidth,
    height=6cm,
    xlabel={Time (seconds)},
    ylabel={Latency (ms)},
    xmin=20, xmax=50,
    ymin=0, ymax=120,
    grid=both,
    legend pos=north east,
]
\addplot[thick, red, mark=none] coordinates {
    (20, 3.2) (22, 3.5) (23, 85) (24, 98) (25, 92) (26, 45) (28, 12) (30, 4.5) (40, 3.2) (50, 3.2)
};
\draw[dashed, blue] (axis cs:25,0) -- (axis cs:25,120) node[pos=0.9, right, font=\tiny] {Action Taken};
\legend{DL Latency}
\end{axis}
\end{tikzpicture}
\caption{Latency recovery during an F2 congestion event. The agent triggers load reduction at $t=25$\,s, restoring latency to baseline within 5 seconds.}
\label{fig:recovery}
\end{figure}

\subsection{LLM Explanation Capability}
Using the Ollama \texttt{llama3.1:8b} backend, the agent generated the following explanation for the F2 event:
\begin{quote}
    \textit{"A latency spike of 85ms was detected alongside a 15\% throughput drop. Since SINR remains stable at 49dB, the root cause is diagnosed as MAC-layer congestion. A 28\% load reduction was executed, which successfully restored the latency to <5ms."}
\end{quote}
Such narratives provide human operators with the "why" behind the autonomous action, which is critical for the adoption of AI in carrier-grade networks.

\subsection{Real-Time Operations Dashboard}
To provide network operators with real-time visibility into the agent's decision-making process, a web-based dashboard was developed. Fig. \ref{fig:dashboard} illustrates the interface during an active F2 fault event.

\begin{figure*}[htbp]
    \centering
    % Replace 'operations_dashboard.png' with your actual screenshot filename
    \includegraphics[width=0.95\textwidth]{operations_dashboard.png}
    \caption{AI5G Real-Time Operations Dashboard showing live telemetry trends, incident timelines, and the LLM-generated post-run narrative for automated remediation events.}
    \label{fig:dashboard}
\end{figure*}

The dashboard integrates live metric trends (latency, throughput, loss, jitter) with the agent's internal state (Mode: ACT/ADVISE). It displays a run summary of incidents and provides a live event stream of the JSON telemetry records processed by the Observer. The bottom section maintains a persistent "Incident Timeline" where each remediation action is logged along with its verification state and the human-readable explanation generated by the LLM.

\subsection{Multi-Seed Evaluation Performance}
To assess the robustness of the NR-aware agent across different stochastic channel realizations, we conducted a 5-run evaluation sweep with seeds $s \in \{7, 8, 9, 10, 11\}$. Each run lasted 120 seconds and included deterministic fault injections at randomized sub-intervals.

As summarized in Table \ref{tab:eval}, the agent processed a total of 17 incidents across the 5 runs. 9 incidents were successfully resolved within the configured cooldown period, while 8 were escalated. The average recovery ratio was measured at 0.5294, resulting in a "FAIR" performance status.

\begin{figure}[htbp]
    \centering
    % Replace 'evaluation_terminal.png' with your actual screenshot filename
    \includegraphics[width=\columnwidth]{evaluation_terminal.png}
    \caption{Terminal output showing the multi-seed evaluation results for seeds 7--11.}
    \label{fig:terminal_results}
\end{figure}

This performance level reflects the challenges of remediating F3 radio degradation faults, which are inherently non-remediable by traffic-layer actions and thus correctly contribute to escalation counts rather than false resolution states.

\begin{table}[htbp]
\centering
\caption{Aggregate Performance Summary (5 Runs, Seeds 7--11)}
\label{tab:eval}
\begin{tabular}{ll}
\toprule
\textbf{Metric} & \textbf{Value} \\
\midrule
Total Incidents & 17 \\
Successfully Resolved & 9 \\
Escalated to Operator & 8 \\
Avg. Recovery Ratio & 0.5294 \\
Evaluation Status & FAIR \\
Total Duration & 600 s \\
\bottomrule
\end{tabular}
\end{table}

\section{Discussion}
The evaluation results highlight a critical property of the proposed system: the high escalation rate for F3 faults. In a 5G NR context, radio link degradation due to UE mobility or shadowing (F3) cannot be resolved by simply restarting the gNB or reducing the traffic load. The agent's ability to identify these as non-remediable and escalate them—rather than attempting infinite remediation loops—is a key feature that prevents service oscillations and maintains operator visibility.

\section{Conclusion}
This project successfully implemented an autonomous, closed-loop fault management system for 5G NR networks. By integrating ns-3 / 5G-LENA with a prioritized AI agent, we demonstrated that incorporating physical layer metrics (SINR, BLER) is essential for accurate fault diagnosis in high-frequency radio environments. The system reliably distinguishes between node-level, traffic-level, and radio-level impairments. The multi-seed evaluation confirmed a consistent "FAIR" status with a 0.529 recovery ratio, primarily driven by the correct escalation of PHY-layer radio link failures.

\section{Future Work}
Future enhancements will focus on several key areas:
\begin{enumerate}
    \item \textbf{Multi-Cell Interference:} Extending the simulation to include adjacent-cell interference to model more complex SINR environments.
    \item \textbf{Deep Reinforcement Learning:} Replacing the rule-based planner with a DQN or PPO agent to optimize remediation for a broader range of handover and beam-management faults.
    \item \textbf{O-RAN Integration:} Packaging the agent as a standard O-RAN xApp/rApp for deployment on production-grade RIC platforms.
    \item \textbf{Explainable AI (XAI):} Enhancing the LLM Explainer to provide SHAP or LIME-based feature importance alongside the narrative descriptions.
    \item \textbf{Energy-Aware Control:} Incorporating gNB power models to support energy-saving modes (Cell DTX/DRX) during low-load periods without compromising fault detection sensitivity.
\end{enumerate}

\end{document}
